\documentclass[a4paper,11pt]{article}
\usepackage[a4paper,margin=2cm]{geometry}
\usepackage[T1]{fontenc}
\usepackage[utf8]{inputenc}
\usepackage[ngerman,english]{babel}
\usepackage{lmodern,microtype,xcolor,parskip,enumitem}
\usepackage[most]{tcolorbox}
\usepackage{titlesec}
\usepackage[hidelinks]{hyperref}
\definecolor{HeaderBlueDark}{RGB}{70,110,170}
\definecolor{RowBlueLight}{RGB}{235,243,255}
\definecolor{RowBlueDark}{RGB}{220,233,250}
\definecolor{SoftGray}{RGB}{90,90,90}
\setlength{\parindent}{0pt}
\setlist[itemize]{leftmargin=1.4em,itemsep=0.35em,topsep=0.35em}
\linespread{1.06}
\titleformat{\section}[block]{\Large\bfseries\color{HeaderBlueDark}}{}{0pt}{}
\titlespacing{\section}{0pt}{1.3em}{0.8em}
\newtcolorbox{exbox}{enhanced,breakable,colback=RowBlueLight,colframe=RowBlueDark,boxrule=0.4pt,arc=1.5mm,left=1.2mm,right=1.2mm,top=1mm,bottom=1mm,title={Beispiele (Examples)},fonttitle=\bfseries,colbacktitle=RowBlueDark,coltitle=black}
\NewDocumentEnvironment{grammarentry}{mm+b}{\begin{tcolorbox}[enhanced,breakable,colback=white,colframe=HeaderBlueDark,boxrule=0.6pt,arc=1.5mm,left=1.5mm,right=1.5mm,top=1mm,bottom=1mm,colbacktitle=HeaderBlueDark,coltitle=white,fonttitle=\bfseries,title={#1\hfill\normalfont\footnotesize #2}]#3\end{tcolorbox}}{}
\newcommand{\GermanText}[1]{\textbf{Deutsch: }#1\par}
\newcommand{\EnglishText}[1]{{\small\color{SoftGray}\textbf{English: }#1\par}}
\newcommand{\ExampleItem}[2]{\item \textbf{#1}\\{\small\color{SoftGray}#2}}
\title{Das Fragewort „warum“ (The question word “warum”)}
\date{}
\begin{document}
\maketitle
\tableofcontents
\newpage
\section{Grundbedeutung und Wortart (Basic meaning and word class)}
\begin{grammarentry}{warum}{Interrogativadverb (kausal)}
\GermanText{\textbf{warum} ist ein kausales Interrogativadverb. Es fragt nach einem Grund: „aus welchem Grund?“.}
\EnglishText{\textbf{warum} is a causal interrogative adverb. It asks for a reason: “for what reason?”.}

\end{grammarentry}
\section{Satzstellung und Verwendung (Word order and use)}
\begin{grammarentry}{Direkte und indirekte Frage}{Fragewortstellung}
\GermanText{In einer direkten W-Frage steht das finite Verb gewöhnlich an zweiter Stelle nach dem Fragewort. In einer indirekten Frage steht es am Ende. In direkten Fragen steht das finite Verb nach \textbf{warum} an zweiter Stelle; in indirekten Fragen steht es am Ende. Auf \textbf{warum} antwortet man oft mit einem \textbf{weil}-Satz.}
\EnglishText{In a direct information question, the finite verb normally follows the question word in second position. In an indirect question, it goes at the end. In direct questions, the finite verb follows \textbf{warum} in second position; in indirect questions, it goes at the end. A \textbf{weil}-clause often answers \textbf{warum}.}
\end{grammarentry}
\begin{exbox}\begin{itemize}
\ExampleItem{Warum lernst du Deutsch? — Weil ich in Österreich lebe.}{Why are you learning German? — Because I live in Austria.}
\ExampleItem{Ich weiß nicht, warum sie abgesagt hat.}{I do not know why she cancelled.}
\end{itemize}\end{exbox}
\section{Zusammenfassung (Summary)}
\begin{grammarentry}{Merksatz}{kurz und wichtig}
\GermanText{Merksatz: \textbf{warum} fragt nach einem Grund.}
\EnglishText{Summary: \textbf{warum} asks for a reason.}
\end{grammarentry}
\end{document}
